\chapter{Pink eye (conjunctivitis)}
\index{Pink eye (conjunctivitis)}

\medskip
\noindent\textit{Educational information; seek professional advice for individual care.}

\section*{What it is}
Conjunctivitis, often called pink eye, is inflammation of the conjunctiva, the thin surface covering the white of the eye and inner eyelid. Common causes include viruses, bacteria, allergies, irritation, chemicals, contact lenses, and foreign bodies.

\section*{Symptoms}
Symptoms may include redness, eyelid swelling, watering, itch, burning, a gritty feeling, mucus or pus, and discomfort. Viral and bacterial forms can spread between people. Eye pain, light sensitivity, persistent blurred vision, or marked redness may indicate a more serious eye problem.

\section*{Care and treatment}
Wash hands, avoid sharing towels, clean discharge gently, and do not rub the eyes. Cold compresses and artificial tears may relieve symptoms. Stop contact lenses until symptoms have resolved and an eye-care professional says it is safe. Antibiotics do not treat viral conjunctivitis; treatment depends on the cause.

\section*{When to seek urgent care}
Seek urgent eye or emergency care for significant eye pain, light sensitivity, blurred vision that does not clear after wiping discharge, intense redness, injury or chemical exposure, a contact-lens-related red eye, severe swelling, or worsening symptoms. Newborns with red or discharging eyes need prompt medical assessment.

\section*{Sources}
\begin{itemize}
\item CDC about pink eye: \url{https://www.cdc.gov/conjunctivitis/about/index.html}
\item CDC treatment guidance: \url{https://www.cdc.gov/conjunctivitis/treatment/index.html}
\end{itemize}
